To understand where LLM/VLM planning actually breaks on a real
robot, we ran a structured investigation of the baseline on
pick-and-place. \rev{We exercised both planner implementations across
all three task families,} and traced every hardware failure back to
its root cause using the baseline's per-run debug artifacts: the raw
model response, the active configuration, and, in VLM mode, the
captured camera frame.
% TODO-VERIFY: "exercised both planner implementations across all
% three task families" — confermare dai log che la copertura è stata
% davvero completa (entrambe le implementazioni x tutte le famiglie di
% task), altrimenti ammorbidire.
The investigation yielded six recurring failure modes, summarized in
Table~\ref{tab:taxonomy}. They fall into three families:
\begin{itemize}
  \item \textbf{Measurement failures}: the plan uses a quantity
    only a sensor can know, but takes the model's guess instead
    (double-counted stacking height, silhouette-edge width sampling).
  \item \textbf{Embodiment failures}: the plan's abstraction of
    the robot or the scene is simpler than the real thing, a point
    where the target has extent, a constant where the gripper's
    geometry varies (phantom top-surface grasp, point-distance target
    matching, width-dependent tool offset).
  \item \textbf{Contact-dynamics failures}: the plan constrains
    final positions but not what bodies do over time (sequential
    multi-object placement collision, still \rev{unresolved}).
\end{itemize}

One finding shapes everything that follows. These failures do not
depend on perception quality. The phantom top-surface grasp appeared
first in the VLM \rev{implementation}, was fixed there, and then
recurred independently in the LLM \rev{implementation}, whose scene
description is exact by construction. The benchmark of
Section~\ref{sec:evaluation} later confirmed the same pattern for the
placement collision: it dominates \emph{both} \rev{implementations}.
The problem is how the model reasons about physical quantities when
turning a decision into numbers, not how the scene was perceived. This
is why the harness of Section~\ref{sec:harness} operates at the
execution level, correcting the numbers a plan carries, and not at the
perception level.

\begin{table*}[t]
\centering
\caption{The six failure modes. ``Physical property required'' marks
what makes each failure impossible to observe in a symbolic or
scripted simulator.}
\label{tab:taxonomy}
\begin{tabularx}{\textwidth}{@{}p{3.1cm}Xp{3.9cm}@{}}
\toprule
\textbf{Failure} & \textbf{Symptom and root cause} & \textbf{Physical property required} \\
\midrule
Phantom top-surface grasp &
  The release lift is set to the object's full height, assuming the
  object was grasped at its top. The executor actually grasps mid-body,
  descending by $\descent = \min(0.5\height, \descent_{\max})$. Net
  effect: a systematic release overshoot of roughly half the object's
  height. Recurred independently in both \rev{implementations}. &
  Real gripper kinematics: mid-body grasping is a deliberate stability
  choice, not a simulator default. \\
\addlinespace
Double-counted stacking height &
  \rev{During camera-based scene grounding, the target's top surface
  was already measured correctly from depth. The planner's stacking
  rule then added the model's separate, unverified height estimate on
  top of it a second time.} &
  A real depth measurement of an already-correct surface height. \\
\addlinespace
Point-distance target matching &
  A \rev{placement} genuinely inside the tray's real footprint, but
  far from its registered center, was scored as a miss by a
  point-distance check. &
  A real target with spatial extent of tens of centimeters, not a
  point. \\
\addlinespace
Silhouette-edge width sampling &
  Grasp width was read from depth exactly on the object's silhouette
  boundary, the single worst location for valid depth due to occlusion
  and parallax. One failed edge sample aborted an otherwise valid
  plan. &
  Real depth noise at object boundaries; absent for a rendered or
  symbolic bounding box. \\
\addlinespace
Width-dependent tool offset &
  The gripper's fingers pivot, so the flange-to-contact distance is a
  function of the commanded opening width, not a constant. A single
  calibrated offset is wrong for any object of a different width. &
  Real finger kinematics, absent from a symbolic action space. \\
\addlinespace
Sequential-placement collision \textbf{(\rev{unresolved})} &
  Zone-aware spacing keeps a zone's \emph{final} release positions from
  overlapping. It does not stop the arm clipping an already-placed
  object en route to the next release, nor a released object settling
  into a neighbor. Dominant residual failure in the 120-trial
  benchmark: all 22 non-empty operator notes describe a collision or a
  fallen object during a multi-object task, in \textbf{both}
  \rev{implementations}. &
  Real multi-body contact between placed objects and the approach
  trajectory, over time. \\
\bottomrule
\end{tabularx}
\end{table*}

\rev{Below we develop the two most instructive resolved cases in
full, then treat the still-unresolved collision failure on its own.}

\subsection{Case study 1: phantom top-surface grasp}
\label{sec:case-phantom}

The executor deliberately grasps objects mid-body rather than at the
top surface. The pick descent is $\descent = \min(0.5\height,
\descent_{\max})$, where $\height$ is the object's height and
$\descent_{\max}$ is a clamp that keeps the gripper's rigid mounting
body away from tall objects. A mid-body grasp is simply more stable
than pinching a top edge.

The paired release action, however, computes the lift height from
the object's \emph{full} height. That computation silently assumes a
top-surface grasp. The assumption is wrong by construction whenever
the pick descended mid-body, which is always. The result is a
systematic overshoot of roughly $0.5\height$: objects are dropped from
a height proportional to half their own size, rather than from
fractions of a centimeter.

We first fixed this in the VLM \rev{implementation}, by deriving the
object's height from measured depth instead of the model's reported
size. \rev{The same failure then recurred, independently, in the LLM
implementation:} its prompt still asked the model to compute the lift
from the scene's declared object height, the exact assumption the
earlier fix had already identified as wrong. Fixing the failure once,
in one \rev{implementation}, did not fix it everywhere it existed. The
error is not a mistake one prompt makes. It is a mismatch between what
the model is asked to assume and what the gripper actually does, and
it therefore needs an architectural fix, not a prompt fix.

\subsection{Case study 2: double-counted stacking height}
\label{sec:case-doublecount}

The scene convention represents each object by a base reference height
and a size, so base plus height reconstructs the top surface. Every
stacking prompt relies on this formula. \rev{In the camera-grounded
pipeline, the scene-grounding step deprojects the target's pixel
location into a depth-measured world coordinate. An early version
wrote that measured value, already the true top surface, into the
scene description as the target's base reference, while leaving the
model's separate, unverified height estimate untouched. The planner
then applied its normal stacking formula and added the guessed height
on top of an already-correct value.} Every stack-on-target release was
inflated by however wrong that guess happened to be.
% TODO-VERIFY: caso riscritto senza citare la modalità ibrida VLM->LLM
% (commento MN). Confermare che la descrizione "camera-grounded
% pipeline / scene-grounding step" corrisponde anche al code path
% della modalità VLM pura; se il bug esisteva SOLO nella modalità
% ibrida, decidere se tenere questo failure mode nella tassonomia
% (toglierlo cambierebbe i conteggi 6/5 in tutto il paper).

The lesson generalizes Case study 1. A sensor-measured quantity must be
inserted \emph{into the representation the planner already expects},
consistently with its conventions, rather than mixed with model
guesses that the planner will keep combining by formula.

\subsection{The \rev{unresolved} case: sequential-placement collision}
\label{sec:case-collision}

Unlike the five other failures, this one is not \rev{resolved} by the
harness. When a task places several objects into the same zone in
sequence, the harness spaces the \emph{final} release positions so
they cannot overlap. That is only the static half of the problem.
Nothing constrains the arm's trajectory while it approaches the
$n$-th release past the $(n-1)$ objects already in the zone. Nothing
constrains a released object's residual motion from displacing a
neighbor afterward, either. Both halves are three-dimensional,
time-dependent contact problems. They cannot be solved by correcting a
single scalar, which is exactly what makes them different in kind from
Case studies 1 and 2.

We diagnosed this failure from the same debug artifacts as the others,
but by the time it was isolated it had already become the dominant
failure mode of the benchmark, in both \rev{implementations}. We
report it as an identified, partially mitigated, \rev{unresolved}
problem; Section~\ref{sec:discussion} returns to what \rev{resolving}
it would require.
